\RequirePackage[l2tabu, orthodox]{nag}
\RequirePackage{silence}
\WarningFilter{nag}{There is no environment ``centering'' }%nag complains because beamer titlepage uses a centering environment
\WarningFilter{nag}{1 complaints in total}
\WarningFilter{ifpdf}{Someone has redefined \pdfoutput}
\WarningFilter{fmtcount}{\ordinal already defined use \FCordinal instead}
\documentclass[english, french]{beamer}
%INSTALL

%avoids a warning
\usepackage[log-declarations=false]{xparse}
\usepackage{fontspec} %font selecting commands
\usepackage{xunicode}
%warn about missing characters
\tracinglostchars=2

%REDAC
\usepackage{booktabs}
\usepackage{calc}

\usepackage{mathtools} %load this before babel!
	\mathtoolsset{showonlyrefs,showmanualtags}

\usepackage{babel}
%suppresses the warning about frenchb not modifying the captions (“—” to “:” in “Figure 1 – Legend”).
	\frenchbsetup{AutoSpacePunctuation=false,SuppressWarning=false}

\usepackage[super]{nth}
\usepackage{listings} %typeset source code listings
	\lstset{language=XML,tabsize=2,literate={"}{{\tt"}}1,captionpos=b}
\usepackage[nolist,smaller,printonlyused]{acronym}%,smaller option produces warnings from relsize in some cases, it seems.
\usepackage[nodayofweek]{datetime}%must be loaded after the babel package
\usepackage{xspace}
\usepackage{hyperref}% option pdfusetitle must be introduced here, not in hypersetup.
%breaklinks makes links on multiple lines into different PDF links to the same target.
%colorlinks (false): Colors the text of links and anchors. The colors chosen depend on the the type of link. In spite of colored boxes, the colored text remains when printing.
%linkcolor=black: this leaves other links in colors, e.g. refs in green, don't print well.
%pdfborder (0 0 1, set to 0 0 0 if colorlinks): width of PDF link border
%hidelinks
\hypersetup{breaklinks,bookmarksopen,colorlinks=true,urlcolor=blue,linkcolor=,hyperfigures=true}
% hyperref doc says: Package bookmark replaces hyperref’s bookmark organization by a new algorithm (...) Therefore I recommend using this package.
\usepackage{bookmark}

% center floats by default, but do not use with float
% \usepackage{floatrow}
% \makeatletter
% \g@addto@macro\@floatboxreset\centering
% \makeatother
\usepackage{ragged2e} %new com­mands \Cen­ter­ing, \RaggedLeft, and \RaggedRight and new en­vi­ron­ments Cen­ter, FlushLeft, and FlushRight, which set ragged text and are eas­ily con­fig­urable to al­low hy­phen­ation (the cor­re­spond­ing com­mands in LaTeX, all of whose names are lower-case, pre­vent hy­phen­ation al­to­gether). 
\usepackage{siunitx} %[expproduct=tighttimes, decimalsymbol=comma]
\sisetup{detect-all}% to detect e.g. when in math mode (use a math font)
\usepackage{braket} %for \Set
\usepackage{natbib}

\usepackage{amsmath,amsthm}
% \usepackage{amsfonts} %not required?
% \usepackage{dsfont} %for what?
%unicode-math overwrites the following commands from the mathtools package: \dblcolon, \coloneqq, \Coloneqq, \eqqcolon. Using the other colon-like commands from mathtools will lead to inconsistencies. Plus, Using \overbracket and \underbracke from mathtools package. Use \Uoverbracket and \Uunderbracke for original unicode-math definition.
%use exclusively \mathbf and choose math bold style below.
\usepackage[warnings-off={mathtools-colon, mathtools-overbracket}, bold-style=ISO]{unicode-math}

\defaultfontfeatures{
	Fractions=On,
	Mapping=tex-text% to turn "--" into dashes, useful for bibtex%%
}
\defaultfontfeatures[\rmfamily, \sffamily]{
	Fractions=On,
	Mapping=% to leave " alone (disable the default mapping tex-text; requires loading the font afterwards?)
}
\newfontfamily\xitsfamily{XITS}
\newfontfamily\texgyretermesfamily{TeX Gyre Termes}
\newfontfamily\texgyreherosfamily{TeX Gyre Heros}
\newfontfamily\lmfamily{Latin Modern Roman}
\setmainfont{Latin Modern Roman}
\setsansfont{Latin Modern Sans}
\setmonofont{Latin Modern Mono}
\defaultfontfeatures{}%disable default font features to avoid warnings with math fonts.
\setmathfont{XITS Math}
\setmathfont[range={\mathcal,\mathbfcal},StylisticSet=1]{XITS Math}

\usepackage{cleveref}% cleveref should go "laster" than hyperref
%GRAPHICS
\usepackage{pgf}
\usepackage{pgfplots}
	\usetikzlibrary{matrix,fit,plotmarks,calc,trees,shapes.geometric,positioning,plothandlers}
\pgfplotsset{compat=1.11}
\usepackage{graphicx}

\graphicspath{{graphics/},{graphics-dm/}}
\DeclareGraphicsExtensions{.pdf}

%HACKING
\usepackage{printlen}
\uselengthunit{mm}
% 	\newlength{\templ}% or LenTemp?
% 	\setlength{\templ}{6 pt}
% 	\printlength{\templ}
\usepackage{etoolbox} %for addtocmd command
\usepackage{scrhack}% load at end. Corrects a bug in float package, which is outdated but might be used by other packages
\usepackage{xltxtra} %somebody said that this is loaded by fontspec, but does not seem correct: if not loaded explicitly, does not appear in the log and \showhyphens is not corrected.

%Beamer-specific
%ADD
\usepackage{appendixnumberbeamer}
%\setbeamersize{text margin left=0.1cm, text margin right=0.1cm} 
\setbeamertemplate{navigation symbols}{}
\usetheme{BrusselsBelgium}
\usefonttheme{professionalfonts}


\newcommand{\R}{ℝ}
\newcommand{\N}{ℕ}
\newcommand{\Z}{ℤ}
\newcommand{\card}[1]{\lvert{#1}\rvert}
\newcommand{\powerset}[1]{\mathscr{P}(#1)}%\mathscr rather than \mathcal: scr is rounder than cal (at least in XITS Math).
\newcommand{\suchthat}{\;\ifnum\currentgrouptype=16 \middle\fi|\;}
%\newcommand{\Rplus}{\reels^+\xspace}

\AtBeginDocument{%
	\renewcommand{\epsilon}{\varepsilon}
% we want straight form of \phi for mathematics, as recommended in UTR #25: Unicode support for mathematics.
%	\renewcommand{\phi}{\varphi}
}

% with amssymb, but I don’t want to use amssymb just for that.
% \newcommand{\restr}[2]{{#1}_{\restriction #2}}
%\newcommand{\restr}[2]{{#1\upharpoonright}_{#2}}
\newcommand{\restr}[2]{{#1|}_{#2}}%sometimes typed out incorrectly within \set.
%\newcommand{\restr}[2]{{#1}_{\vert #2}}%\vert errors when used within \Set and is typed out incorrectly within \set.
\DeclareMathOperator*{\argmax}{arg\,max}
\DeclareMathOperator*{\argmin}{arg\,min}


%ARG TH
\newcommand{\AF}{\mathcal{AF}}
\newcommand{\labelling}{\mathcal{L}}
\newcommand{\labin}{\textbf{in}\xspace}
\newcommand{\labout}{\textbf{out}}
\newcommand{\labund}{\textbf{undec}\xspace}
\newcommand{\nonemptyor}[2]{\ifthenelse{\equal{#1}{}}{#2}{#1}}
\newcommand{\gextlab}[2][]{
	\labelling{\mathcal{GE}}_{(#2, \nonemptyor{#1}{\ibeatsr{#2}})}
}
\newcommand{\allargs}{A^*}
\newcommand{\args}{A}
\newcommand{\ar}{a}

%MCDA+Arg
\newcommand{\dm}{d}
\newcommand{\ileadsto}{\leadsto}
\newcommand{\mleadsto}[1][\eta]{\leadsto_{#1}}
\newcommand{\ibeats}{\vartriangleright}
\newcommand{\mbeats}[1][\eta]{\vartriangleright_{#1}}


%MISC
\newcommand{\lequiv}{\Vvdash}
\newcommand{\weightst}{W^{\,t}}

%MCDA classical
\newcommand{\crits}{\mathcal{J}}
\newcommand{\altspace}{\mathbb{A}}
\newcommand{\alts}{A}

%Sorting
\newcommand{\cats}{\mathcal{C}}
\newcommand{\catssubsets}{2^\cats}
\newcommand{\catgg}{\vartriangleright}
\newcommand{\catll}{\vartriangleleft}
\newcommand{\catleq}{\trianglelefteq}
\newcommand{\catgeq}{\trianglerighteq}
\newcommand{\alttoc}[2][x]{(#1 \xrightarrow{} #2)}
\newcommand{\alttocat}[3]{(#2 \xrightarrow{#1} #3)}
\newcommand{\alttoI}{(x \xrightarrow{} \left[\underline{C_x}, \overline{C_x}\right])}
\newcommand{\alttocatdm}[3][t]{\left(#2 \thinspace \raisebox{-3pt}{$\xrightarrow{#1}$}\thinspace #3\right)}
\newcommand{\alttocatatleast}[2]{\left(#1 \thinspace \raisebox{-3pt}{$\xrightarrow[]{≥}$}\thinspace #2\right)}
\newcommand{\alttocatatmost}[2]{\left(#1 \thinspace \raisebox{-3pt}{$\xrightarrow[]{≤}$}\thinspace #2\right)}

\newcommand{\source}{\scriptsize}
\newcommand{\commentOC}[1]{{\selectlanguage{french}{\todo{OC : #1}}}}
%Or: \todo[color=green!40]

%this probably requires outdated float package, see doc KomaScript for an alternative.
% \newfloat{program}{t}{lop}
% \floatname{program}{PM}

%\crefname{axiom}{axiom}{axioms}%might be needed for workaround bug in cref when defining new theorems?

%\ifdefined\theorem\else
%\newtheorem{theorem}{\iflanguage{english}{Theorem}{Théorème}}
%\fi

%which line breaks are chosen: accept worse lines, therefore reducing risk of overfull lines. Default = 200
\tolerance=2000
%accept overfull hbox up to...
\hfuzz=2cm
%reduces verbosity about the bad line breaks
\hbadness 5000
%sloppy sets tolerance to 9999
\apptocmd{\sloppy}{\hbadness 10000\relax}{}{}

% WRITING
%\newcommand{\ie}{i.e.\@\xspace}%to try
%\newcommand{\eg}{e.g.\@\xspace}
%\newcommand{\etal}{et al.\@\xspace}
\newcommand{\ie}{i.e.\ }
\newcommand{\eg}{e.g.\ }
\newcommand{\mkkOK}{\checkmark}%\color{green}{\checkmark}
\newcommand{\mkkREQ}{\ding{53}}%requires pifont?%\color{green}{\checkmark}
\newcommand{\mkkNO}{}%\text{\color{red}{\textsf{X}}}

\makeatletter
\newcommand{\boldor}[2]{%
	\ifnum\strcmp{\f@series}{bx}=\z@
		#1%
	\else
		#2%
	\fi
}
\newcommand{\textstyleElProm}[1]{\boldor{\MakeUppercase{#1}}{\textsc{#1}}}
\makeatother
\newcommand{\electre}{\textstyleElProm{Électre}\xspace}
\newcommand{\electreIv}{\textstyleElProm{Électre Iv}\xspace}
\newcommand{\electreIV}{\textstyleElProm{Électre IV}\xspace}
\newcommand{\electreIII}{\textstyleElProm{Électre III}\xspace}
\newcommand{\electreTRI}{\textstyleElProm{Électre Tri}\xspace}
% \newcommand{\utadis}{\texorpdfstring{\textstyleElProm{utadis}\xspace}{UTADIS}}
% \newcommand{\utadisI}{\texorpdfstring{\textstyleElProm{utadis i}\xspace}{UTADIS I}}

%TODO
% \newcommand{\textstyleElProm}[1]{{\rmfamily\textsc{#1}}} 


\newlength{\GraphsNodeSep}
\setlength{\GraphsNodeSep}{7mm}

% MCDA Drawing Sorting
\newlength{\MCDSCatHeight}
\setlength{\MCDSCatHeight}{6mm}
\newlength{\MCDSAltHeight}
\setlength{\MCDSAltHeight}{4mm}
%separation between two vertical alts
\newlength{\MCDSAltSep}
\setlength{\MCDSAltSep}{2mm}
\newlength{\MCDSCatWidth}
\setlength{\MCDSCatWidth}{3cm}
\newlength{\MCDSEvalRowHeight}
\setlength{\MCDSEvalRowHeight}{6mm}
\newlength{\MCDSAltsToCatsSep}
\setlength{\MCDSAltsToCatsSep}{1.5cm}
\newcounter{MCDSNbAlts}
\newcounter{MCDSNbCats}
\newlength{\MCDSArrowDownOffset}
\setlength{\MCDSArrowDownOffset}{0mm}

\tikzset{/Graphs/dot/.style={
	shape=circle, fill=black, inner sep=0, minimum size=1mm
}}
\tikzset{/MC/D/S/alt/.style={
	shape=rectangle, draw=black, inner sep=0, minimum height=\MCDSAltHeight, minimum width=2.5cm, anchor=north east
}}
\tikzset{MC/D/S/pref/.style={
	shape=ellipse, draw=gray, thick
}}
\tikzset{/MC/D/S/cat/.style={
	shape=rectangle, draw=black, inner sep=0, minimum height=\MCDSCatHeight, minimum width=\MCDSCatWidth, anchor=north west
}}
\tikzset{/MC/D/S/evals matrix/.style={
	matrix, row sep=-\pgflinewidth, column sep=-\pgflinewidth, nodes={shape=rectangle, draw=black, inner sep=0mm, text depth=0.5ex, text height=1em, minimum height=\MCDSEvalRowHeight, minimum width=12mm}, nodes in empty cells, matrix of nodes, inner sep=0mm, outer sep=0mm, row 1/.style={nodes={draw=none, minimum height=0em, text height=, inner ysep=1mm}}
}}

% Beliefs
\tikzset{/Beliefs/D/S/attacker/.style={
	shape=rectangle, draw, minimum size=8mm
}}
\tikzset{/Beliefs/D/S/supporter/.style={
	shape=circle, draw
}}

\newcommand{\tikzmark}[1]{%
	\tikz[overlay, remember picture, baseline=(#1.base)] \node (#1) {};%
}


\begin{acronym}
\acro{AMCD}{Aide Multicritère à la Décision}
\acro{ASA}{Argument Strength Assessment}
\acro{DA}{Decision Analysis}
\acro{DM}{Decision Maker}
\acro{DPr}{Deliberated Preferences}
\acro{DRSA}{Dominance-based Rough Set Approach}
\acro{DSS}{Decision Support Systems}
\acrodefplural{DSS}{Decision Support Systems}
% \newacroplural{DSS}[DSSes]{Decision Support Systems}
\acro{EJOR}{European Journal of Operational Research}
\acro{LNCS}{Lecture Notes in Computer Science}
\acro{MCDA}{Multicriteria Decision Aid}
\acro{MIP}{Mixed Integer Program}
\acro{NCSM}{Non Compensatory Sorting Model}
\acro{PL}{Programme Linéaire}
\acro{PLNE}{Programme Linéaire en Nombres Entiers}
\acro{PM}{Programme Mathématique}
\acro{MP}{Mathematical Program}
\acro{MIP}{Mixed Integer Program}
% \newacroplural{PM}{Programmes Mathématiques}
%acrodefplural since version 1.35, my debian has \ProvidesPackage{acronym}[2009/01/25, v1.34, Support for acronyms (Tobias Oetiker)]
\acrodefplural{PM}{Programmes Mathématiques}
\acro{PMML}{Predictive Model Markup Language}
\acro{RESS}{Reliability Engineering \& System Safety}
\acro{SMAA}{Stochastic Multicriteria Acceptability Analysis}
\acro{URPDM}{Uncertainty and Robustness in Planning and Decision Making}
\acro{XML}{Extensible Markup Language}
\end{acronym}


\title{Conception d’applications internet}
\subtitle{EJB}
\subject{Java EE}
\keywords{stateful, stateless, session beans, SFSB, SLSB, injection, CDI, resources}
\author{Olivier Cailloux}
\institute[LAMSADE]{LAMSADE, Université Paris-Dauphine}
\date{\formatdate{14}{12}{2015}}

\begin{document}
\bibliographystyle{apalike}

\begin{frame}[plain]
	\tikz[remember picture,overlay]{
		\path (current page.south west) node[anchor=south west, inner sep=0] {
			\includegraphics[height=1cm]{LAMSADE95.jpg}
		};
		\path (current page.south) ++ (0, 1mm) node[anchor=south, inner sep=0] {
			\includegraphics[height=9mm]{Dauphine.jpg}
		};
		\path (current page.south east) node[anchor=south east, inner sep=0] {
			\includegraphics[height=1cm]{PSL.png}
		};
	}
   \titlepage
\end{frame}
\addtocounter{framenumber}{-1}

\section{Annotations}
\subsection{Présentation}
\begin{frame}
	\frametitle{Annotations : présentation}
	\begin{itemize}
		\item Une forme de métadonnée
		\item Détection d’erreurs supplémentaires
		\item Configuration du compilateur
		\item Génération de code (XML, etc.)
		\item Traitement runtime possible
	\end{itemize}
\end{frame}

\subsection{Types}
\begin{frame}[fragile]
	\frametitle{Types annotation}
	\begin{lstlisting}
@interface NameAndNumber {
  String name();
  int number() default 1;
}
	\end{lstlisting}
	\begin{itemize}
		\item \emph{Type} annotation : sorte d’interface ; type doté d’\emph{élements}
		\item Élément : type, identifiant, valeur par défaut
		\item Si pas d’élément, l’annotation est un \emph{Marqueur}
		\item Si un seul élément, élément souvent nommé \texttt{value}
		\item Prédéfinies ou définies par l’utilisateur
	\end{itemize}
\end{frame}

\subsection{Annotations}
\begin{frame}
	\frametitle{Annotations}
	\texttt{@Deprecated} ; \texttt{@SuppressWarnings(value = "unchecked")} ; \texttt{@SuppressWarnings("unchecked")} ; \texttt{@NameAndNumber(name = "Benjamin Franklin", number = 3)}
	\begin{itemize}
		\item Annotation : nom d’un type annotation et valeurs pour ses éléments
		\item Si seulement une valeur : utilisée pour \texttt{value}
		\item Appliquées sur des déclarations (classes, champs, méthodes, usages de types) ; multiples annotations pour une déclaration permises
	\end{itemize}
	Exemple d’usage : \texttt{@NonNull String str;}
\end{frame}

\section[Dernier épisode]{Durant le dernier épisode…}
\subsection{Était à faire}
\begin{frame}
	\frametitle{Était à faire}
	\begin{itemize}
		\item Groupe sur  \href{https://mycourse.dauphine.fr/webapps/blackboard/execute/courseMain?course_id=_34753_1}{MyCourse} \& serveur git : cf. dernière annonce
		\item Exercices servlets : attention au type de ressource renvoyée ; à la gestion des erreurs…
		\item Git : gestion d’un merge avec conflit pour chaque membre de l’équipe
	\end{itemize}
\end{frame}

\section{EJB}
\subsection{Module EJB}
\begin{frame}
	\frametitle{Module EJB}
	\begin{block}{Module EJB}
		\begin{itemize}
			\item Fichier \texttt{.jar}
			\item Fichiers \texttt{.class} EJB et autres
			\item \texttt{META-INF/ejb-jar.xml} : descripteur pour conteneur EJB (attributs de transaction, sécurité, …)
			\item \texttt{META-INF/glassfish-ejb-jar.xml} : descripteur pour glassfish
			\item \texttt{META-INF/MANIFEST.MF}
		\end{itemize}
	\end{block}
\end{frame}

\subsection{Sortes d’EJB}
\begin{frame}
	\frametitle{Sortes d’EJB}
	\begin{itemize}
		\item Session bean ; message-driven bean {\tiny (et entity bean, facultatif)} {\tiny (diapos valables pour session beans)}
		\item Session bean : \href{http://docs.oracle.com/javaee/7/api/javax/ejb/Stateless.html}{\texttt{@Stateless}} (SLSB), \href{http://docs.oracle.com/javaee/7/api/javax/ejb/Stateful.html}{\texttt{@Stateful}} (SFSB) ou \href{http://docs.oracle.com/javaee/7/api/javax/ejb/Singleton.html}{\texttt{@Singleton}}
		\item Stateless : conteneur instancie et gère un pool de beans
		\item Stateful et Singleton : ont un état
		\item Singleton : conteneur instancie et s’assure de l’unicité
		\item Stateful : norme EJB ne gère pas son cycle de vie (à gérer par le client {\tiny sauf si CDI})
		\item Stateful peut être \emph{passivated} : implémenter \href{https://docs.oracle.com/javase/8/docs/api/java/io/Serializable.html}{Serializable} {\tiny (facultatif)}
		\item Singleton : conteneur s’assure de l’unicité
		\item Tous types : thread-safe {\tiny (mais attention aux objects accessibles depuis le bean)}
	\end{itemize}
\end{frame}

\begin{frame}
	\frametitle{Types d’exposition}
	\begin{itemize}
		\item Le conteneur fournit une vue client pour donner accès
		\item Pattern ? \pause Proxy ! \pause
		\item Le développeur déclare comment l’EJB sera accessible
		\begin{itemize}
			\item No-interface view : accès direct et local à l’EJB
			\item Local view : accès local via une interface \og{}business\fg{}
			\item Remote view : accès via une interface \og{}business\fg{}
		\end{itemize}
		\item {\tiny Ou accès via vue décrite en WSDL via endpoint JAX-WS, pour session bean stateless ou singleton}
		\item Accès local : client et EJB dans même application Java EE {\tiny (peut être relaché hors standard)}
		\item Inconvénient de l’accès Remote ? \pause Sérialisation !\pause
		\item Recommandations. No-interface view : marquer le bean \href{http://docs.oracle.com/javaee/7/api/javax/ejb/LocalBean.html}{\texttt{@LocalBean}}. Business interface : l’implémenter et la marquer \href{http://docs.oracle.com/javaee/7/api/javax/ejb/Local.html}{\texttt{@Local}} ou \href{http://docs.oracle.com/javaee/7/api/javax/ejb/Remote.html}{\texttt{@Remote}}.
	\end{itemize}
\end{frame}

\begin{frame}
	\frametitle{Convention de nommage}
	\begin{itemize}
		\item Business interface : truc
		\item Classe (et nom du bean) : trucBean
	\end{itemize}
\end{frame}

\subsection{Invocation}
\begin{frame}
	\frametitle{Invocation}
	\begin{block}{SLSB et Singleton session beans}
		\begin{itemize}
			\item Injecter EJB via DI : annoter la référence \href{http://docs.oracle.com/javaee/7/api/javax/ejb/EJB.html}{\texttt{@EJB}} {\tiny (v. aussi @Inject via CDI)}
			\item À utiliser pour les EJB dont le conteneur gère le cycle de vie !
		\end{itemize}
	\end{block}
	Conteneur rend EJBs accessible via JNDI : \texttt{java:global[/<app>]/<module>/<bean>[!<interface>]} ; \texttt{java:app/<module>/<bean>[!<interface>]} ; \texttt{java:module/<bean>[!<interface>]}
	\begin{block}{SFSB}
		\begin{itemize}
			\item Création explicite via JNDI {\tiny (sauf gestion via CDI)}
			\item Destruction explicite via méthode EJB annotée \href{http://docs.oracle.com/javaee/7/api/javax/ejb/Remove.html}{@Remove}
		\end{itemize}
	\end{block}
\end{frame}

\subsection{JNDI}
\begin{frame}
	\frametitle{JNDI : aperçu}
	\begin{block}{JNDI}
		\begin{itemize}
			\item JNDI : Java Naming and Directory Interface (Java SE)
			\item API pour accéder à des objets via leur nom
			\item Et à des \og{}naming contexts\fg{}
			\item Indépendant du stockage du répertoire
		\end{itemize}
	\end{block}
	\begin{block}{Accès}
		\begin{itemize}
			\item Créer un \href{http://docs.oracle.com/javase/8/docs/api/javax/naming/InitialContext.html}{\texttt{InitialContext}}
			\item Utiliser méthodes de \href{http://docs.oracle.com/javase/8/docs/api/javax/naming/Context.html}{\texttt{Context}} t.q. \texttt{lookup(String name)}
			\item Accès en lecture uniquement dans cas EJB
		\end{itemize}
	\end{block}
\end{frame}

\subsection{Environnement}
\begin{frame}
	\frametitle{Environnement}
	\begin{itemize}
		\item L’EJB a un environnement
		\item Accessible via JNDI dans le naming context \texttt{java:comp/env}
		\item \emph{Ressource} : référence dans env. d’un composant Java EE
		\item Permet la paramétrisation de l’EJB (exemple : valeur entière)
		\item Depuis EJB, accéder via JNDI
		\item Ou utiliser CDI et \href{http://docs.oracle.com/javaee/7/api/javax/annotation/Resource.html}{\texttt{@Resource}}
		\item Déclaration : où ? \pause Dans descripteur standard EJB.\pause
		\item Resource non injectée si pas de valeur dans le descripteur : utilise la valeur donnée dans le code
	\end{itemize}
\end{frame}

\begin{frame}[fragile]
	\frametitle{Déclarer et accéder à une ressource}
	Exemple de déclaration d’une ressource (de type \og{}simple\fg{}) \texttt{java:comp/env/ejb.MySesEjbBean/zeNumber} via \texttt{ejb-jar.xml}
	\lstinputlisting[firstline=6, lastline=15, basicstyle=\scriptsize]{ejb-jar-zeNumber.xml}
	Accès à la ressource
	\lstinputlisting[firstline=12, lastline=15, basicstyle=\scriptsize]{MySesEjbBean.java}
\end{frame}

\subsection{Responsabilités}
\begin{frame}
	\frametitle{Responsabilités}
	\begin{itemize}
		\item L’EJB doit avoir des responsabilités business
		\item À usage interne {\tiny (typiquement)}
		\item EJB {\tiny (et data model)} fournissent le modèle et les services
		\item Les servlets {\tiny au sens large} fournissent une vue
		\item \emph{Contrat} clair requis… comme toujours !
		\item Contrat : devoirs de l’utilisateur ; droits attendus en échange
		\item À indiquer dans la javadoc
		\item Utiliser des copies défensives ! Pourquoi ? \pause Toujours bien ; important si EJB devient Remote
	\end{itemize}
\end{frame}

\subsection{Références}
\begin{frame}
	\frametitle{Références}
	\begin{itemize}
		\item The Java EE Tutorial: \href{http://docs.oracle.com/javaee/7/tutorial/partentbeans.htm}{Enterprise Beans}
		\item \href{https://www.jcp.org/en/jsr/detail?id=345}{JSR 345} (EJB 3.2) (\href{http://download.oracle.com/otn-pub/jcp/ejb-3_2-fr-eval-spec/ejb-3_2-core-fr-spec.pdf}{direct})
		\item \href{https://www.jcp.org/en/jsr/detail?id=342}{JSR 342} (Java EE 7)
		\item \href{http://jcp.org/en/jsr/detail?id=346}{JSR 346} (Context and Dependency Injection 1.1 et 1.2) (\href{http://download.oracle.com/otn-pub/jcp/cdi-1_2-mrel-eval-spec/cdi-1.2.pdf}{direct}).
		\item \href{http://docs.oracle.com/javase/8/docs/technotes/guides/jndi/}{JNDI doc}
	\end{itemize}
\end{frame}

\section{Servlets}
\subsection{Session HTTP}
\begin{frame}
	\frametitle{Servlet et session HTTP}
	\begin{itemize}
		\item Un servlet s’exécute dans un contexte
		\item Entre autres, une session HTTP
		\item Objet \href{http://docs.oracle.com/javaee/7/api/javax/servlet/http/HttpSession.html}{\texttt{HttpSession}} rendu disponible par conteneur ({\tiny injection ou } accès via paramètres servlet)
		\item Le serveur tente de tracker la session via un cookie {\tiny par exemple}
		\item Le serveur fait expirer le cookie après un temps configurable
		\item Un bon site fonctionne même sans cookies !
		\item Classe implémente \href{http://docs.oracle.com/javaee/7/api/javax/servlet/http/HttpSessionListener.html}{\texttt{HttpSessionListener}} pour recevoir notifications \texttt{sessionCreated}, \texttt{sessionDestroyed}.  L’annoter \href{http://docs.oracle.com/javaee/7/api/javax/servlet/annotation/WebListener.html}{\texttt{@WebListener}}. (Pattern ? \pause Observateur !)
	\end{itemize}
\end{frame}

\section{Sérialisation}
\begin{frame}
	\frametitle{Sérialisation}
	\begin{itemize}
		\item Permet à un objet d’être écrit sur disque
		\item Implémenter \href{http://docs.oracle.com/javase/8/docs/api/java/io/Serializable.html}{Serializable}
		\item Interface marqueur. Exemple de quel principe ? \pause Contrat !
		\item Indique engagement du développeur \pause
		\item Documenter les champs sérialisés : \href{http://docs.oracle.com/javase/1.5.0/docs/tooldocs/windows/javadoc.html@serial}{\texttt{@serial}} (javadoc tag)
		\item Marquer les champs non-sérialisés \texttt{transient} (éphémère)
		\item Enfants doivent aussi être sérialisables
		\item {\tiny Cas général : déclarer un \texttt{serialVersionUID} ; peu important pour les objets non transmis t.q. SFSB}
		\item Voir \href{https://docs.oracle.com/javase/8/docs/technotes/guides/serialization/index.html}{Java Object Serialization} (FAQ utile)
	\end{itemize}
\end{frame}

\section{Exercices}
\subsection{EJB}
\begin{frame}[allowframebreaks]
	\frametitle{Exercices : EJB}
	\begin{itemize}
		\item Programmer un SLSB qui implémente une méthode \texttt{add(int, int): int}. Il s’occupera de la partie algorithmique de votre premier servlet.
		\item Votre EJB logge (au niveau \texttt{info}) les valeurs reçues et renvoyées
		\item Si le résultat est nul, votre EJB logge ce fait (en plus) au niveau \texttt{warn}
		\item Programmer un servlet qui accède à l’EJB pour additionner deux paramètres fournis par le client web
		\item Déployer et tester l’application résultante ; vérifier les logs
		\item Ajouter une ressource à votre EJB, valant 1 par défaut, qui est ajoutée au résultat du calcul avant renvoi à l’appelant. Tester le servlet.
		\item Exporter cette ressource, prendre son chapeau de responsable du déploiement, et modifier la valeur par défaut de la ressource, sans toucher au code source donc. Vérifier l’effet.
		\item Programmer un SFSB qui incrémente un compteur à l’aide d’une valeur passée en paramètre (reçue via votre servlet)
		\item Utiliser votre SFSB à la place de votre SLSB dans votre servlet. Durant toute la session HTTP utilisateur, l’utilisateur voit donc le compteur augmenter avec les valeurs qu’il fournit en paramètre.
		\item Programmer la destruction de votre SFSB à la destruction de la session. Que se passe-t-il si vous ne le faites pas ?
		\item Programmer l’expiration de la session HTTP après un délai court, via votre servlet (indice : nous avons déjà rencontré l’objet qui permet de le faire)
		\item Vérifier (via les logs) que votre SFSB est bien créé et détruit à l’ouverture et à l’expiration des sessions HTTP. Utiliser pour ce faire différents onglets ou fenêtres de votre navigateur et la fonction de suppression des cookies.
	\end{itemize}
\end{frame}

\section{Rappel méthodologique}
\subsection{Cycles courts}
\begin{frame}
	\frametitle{Développement agile \& cycles courts}
	\begin{itemize}
		\item Développement agile : release often
		\item Cycles courts : applications complètes !
		\item Réfléchir aux fonctionnalités prioritaires
		\item Attention aux mélanges cycle en V et agile !
	\end{itemize}
\end{frame}

\section{À faire}
\begin{frame}
	\frametitle{À faire}
	\begin{block}{Avant le 21 janvier}
		\begin{itemize}
			\item[+\footnote{Le + indique que cet aspect intervient dans la note}] Chaque binôme : \emph{au moins} un servlet utile à votre projet
			\item[+] Chaque binôme : \emph{au moins} un EJB utile à votre projet
			\item[+] Au moins un push par personne dans le projet
		\end{itemize}
	\end{block}
	Chaque commit doit indiquer clairement le ou les auteurs
	\begin{itemize}
		\item Utilisez \texttt{user.name}
		\item Indiquez dans le commentaire du commit le binôme éventuel
	\end{itemize}
\end{frame}

\appendix
\AtBeginSection{
}
\section{Licence}
\begin{frame}
	\frametitle{Licence}
	Cette présentation, et le code LaTeX associé, sont sous \href{http://opensource.org/licenses/MIT}{licence MIT}. Vous êtes libres de réutiliser des éléments de cette présentation, sous réserve de citer l’auteur.
	
	Le travail réutilisé est à attribuer à \href{http://www.lamsade.dauphine.fr/~ocailloux/}{Olivier Cailloux}, Université Paris-Dauphine.
	
	\small{(Ceci ne couvre pas les images incluses dans ce document, puisque je n’en suis généralement pas l’auteur.)}
\end{frame}
\end{document}
\begin{frame}
	\frametitle{}
	\begin{itemize}
		\item 
	\end{itemize}
	\begin{block}{}
		
	\end{block}
\end{frame}

\section{Bibliographie}
\begin{frame}[allowframebreaks]
	\frametitle{Bibliographie}
	\def\newblock{\hskip .11em plus .33em minus .07em}
% 	\bibliography{zotero}
\end{frame}

\section{Autres}
\begin{frame}
	\frametitle{}
	\begin{itemize}
		\item 
	\end{itemize}
	\begin{block}{}
		
	\end{block}
\end{frame}
\end{document}
